% problem_id: S001-lemma-huge-diagram-ringed-topoi
% source_id: S001 (Stacks Project)
% source_file: defos.tex
% statement_locator: defos.tex lines 3905--3978
% proof_locator: defos.tex lines 3980--4133
% env: lemma
% id_note: label 在本来源内唯一
% pairing: tight (verified)
% status: complete (statement + author proof verbatim + reason + locators)
%
\begin{lemma}
\label{lemma-huge-diagram-ringed-topoi}
Assume given a commutative diagram of morphisms ringed topoi
\begin{equation}
\label{equation-huge-1-ringed-topoi}
\vcenter{
\xymatrix{
& (\Sh(\mathcal{C}_2), \mathcal{O}_2) \ar[r]_{i_2} \ar[d]_{f_2} \ar[ddl]_g &
(\Sh(\mathcal{C}'_2), \mathcal{O}'_2) \ar[d]^{f'_2} \\
&
(\Sh(\mathcal{B}_2), \mathcal{O}_{\mathcal{B}_2}) \ar[r]^{t_2} \ar[ddl]|\hole &
(\Sh(\mathcal{B}'_2), \mathcal{O}_{\mathcal{B}'_2}) \ar[ddl] \\
(\Sh(\mathcal{C}_1), \mathcal{O}_1) \ar[r]_{i_1} \ar[d]_{f_1} &
(\Sh(\mathcal{C}'_1), \mathcal{O}'_1) \ar[d]^{f'_1} \\
(\Sh(\mathcal{B}_1), \mathcal{O}_{\mathcal{B}_1}) \ar[r]^{t_1} &
(\Sh(\mathcal{B}'_1), \mathcal{O}_{\mathcal{B}'_1})
}
}
\end{equation}
whose horizontal arrows are first order thickenings. Set
$\mathcal{G}_j = \Ker(i_j^\sharp)$ and assume given a
map of $g^{-1}\mathcal{O}_1$-modules
$\nu : g^{-1}\mathcal{G}_1 \to \mathcal{G}_2$
giving rise to the commutative diagram
\begin{equation}
\label{equation-huge-2-ringed-topoi}
\vcenter{
\xymatrix{
& 0 \ar[r] & \mathcal{G}_2 \ar[r] &
\mathcal{O}'_2 \ar[r] &
\mathcal{O}_2 \ar[r] & 0 \\
& 0 \ar[r]|\hole &
f_2^{-1}\mathcal{J}_2 \ar[u]_{c_2} \ar[r] &
f_2^{-1}\mathcal{O}_{\mathcal{B}'_2} \ar[u] \ar[r]|\hole &
f_2^{-1}\mathcal{O}_{\mathcal{B}_2} \ar[u] \ar[r] & 0 \\
0 \ar[r] &
\mathcal{G}_1 \ar[ruu] \ar[r] &
\mathcal{O}'_1 \ar[r] &
\mathcal{O}_1 \ar[ruu] \ar[r] & 0 \\
0 \ar[r] &
f_1^{-1}\mathcal{J}_1 \ar[ruu]|\hole \ar[u]^{c_1} \ar[r] &
f_1^{-1}\mathcal{O}_{\mathcal{B}'_1} \ar[ruu]|\hole \ar[u] \ar[r] &
f_1^{-1}\mathcal{O}_{\mathcal{B}_1} \ar[ruu]|\hole \ar[u] \ar[r] & 0
}
}
\end{equation}
with front and back solutions to (\ref{equation-to-solve-ringed-topoi}).
(The north-north-west arrows are maps on $\mathcal{C}_2$ after applying
$g^{-1}$ to the source.)
\begin{enumerate}
\item There exist a canonical element in
$\Ext^1_{\mathcal{O}_2}(
Lg^*\NL_{\mathcal{O}_1/\mathcal{O}_{\mathcal{B}_1}}, \mathcal{G}_2)$
whose vanishing is a necessary and sufficient condition for the existence
of a morphism of ringed topoi
$(\Sh(\mathcal{C}'_2), \mathcal{O}'_2) \to
(\Sh(\mathcal{C}'_1), \mathcal{O}'_1)$ fitting into
(\ref{equation-huge-1-ringed-topoi}) compatibly with $\nu$.
\item If there exists a morphism
$(\Sh(\mathcal{C}'_2), \mathcal{O}'_2) \to
(\Sh(\mathcal{C}'_1), \mathcal{O}'_1)$
fitting into
(\ref{equation-huge-1-ringed-topoi}) compatibly with $\nu$ the set
of all such morphisms is a principal homogeneous space under
$$
\Hom_{\mathcal{O}_1}(
\Omega_{\mathcal{O}_1/\mathcal{O}_{\mathcal{B}_1}}, g_*\mathcal{G}_2) =
\Hom_{\mathcal{O}_2}(
g^*\Omega_{\mathcal{O}_1/\mathcal{O}_{\mathcal{B}_1}}, \mathcal{G}_2) =
\Ext^0_{\mathcal{O}_2}(
Lg^*\NL_{\mathcal{O}_1/\mathcal{O}_{\mathcal{B}_1}}, \mathcal{G}_2).
$$
\end{enumerate}
\end{lemma}

\begin{proof}
The proof of this lemma is identical to the proof of
Lemma \ref{lemma-huge-diagram-ringed-spaces}.
We urge the reader to read that proof instead of this one.
We will identify the underlying topoi for every
thickening in sight (we have already used this convention
in the statement). The equalities in the last statement of the
lemma are immediate from the definitions. Thus we will work with the groups
$\Ext^k_{\mathcal{O}_2}(
Lg^*\NL_{\mathcal{O}_1/\mathcal{O}_{\mathcal{B}_1}}, \mathcal{G}_2)$,
$k = 0, 1$ in the rest of the proof. We first argue that we can reduce
to the case where the underlying topos of all ringed topoi in the lemma
is the same.

\medskip\noindent
To do this, observe that
$g^{-1}\NL_{\mathcal{O}_1/\mathcal{O}_{\mathcal{B}_1}}$ is equal to the naive
cotangent complex of the homomorphism of sheaves of rings
$g^{-1}f_1^{-1}\mathcal{O}_{\mathcal{B}_1} \to g^{-1}\mathcal{O}_1$, see
Modules on Sites, Lemma \ref{sites-modules-lemma-pullback-differentials}.
Moreover, the degree $0$ term of
$\NL_{\mathcal{O}_1/\mathcal{O}_{\mathcal{B}_1}}$ is a flat
$\mathcal{O}_1$-module, hence the canonical map
$$
Lg^*\NL_{\mathcal{O}_1/\mathcal{O}_{\mathcal{B}_1}}
\longrightarrow
g^{-1}\NL_{\mathcal{O}_1/\mathcal{O}_{\mathcal{B}_1}}
\otimes_{g^{-1}\mathcal{O}_1} \mathcal{O}_2
$$
induces an isomorphism on cohomology sheaves in degrees $0$ and $-1$.
Thus we may replace the Ext groups of the lemma with
$$
\Ext^k_{g^{-1}\mathcal{O}_1}(
g^{-1}\NL_{\mathcal{O}_1/\mathcal{O}_{\mathcal{B}_1}}, \mathcal{G}_2) =
\Ext^k_{g^{-1}\mathcal{O}_1}(
\NL_{g^{-1}\mathcal{O}_1/g^{-1}f_1^{-1}\mathcal{O}_{\mathcal{B}_1}},
\mathcal{G}_2)
$$
The set of morphism of ringed topoi
$(\Sh(\mathcal{C}'_2), \mathcal{O}'_2) \to
(\Sh(\mathcal{C}'_1), \mathcal{O}'_1)$ fitting into
(\ref{equation-huge-1-ringed-topoi}) compatibly with $\nu$ is in
one-to-one bijection with the set of homomorphisms of
$g^{-1}f_1^{-1}\mathcal{O}_{\mathcal{B}'_1}$-algebras
$g^{-1}\mathcal{O}'_1 \to \mathcal{O}'_2$ which are compatible with
$f^\sharp$ and $\nu$. In this way we see that we may assume we have a
diagram (\ref{equation-huge-2-ringed-topoi}) of sheaves on a site
$\mathcal{C}$ (with $f_1 = f_2 = \text{id}$ on underlying topoi)
and we are looking to find a homomorphism of sheaves of rings
$\mathcal{O}'_1 \to \mathcal{O}'_2$ fitting into it.

\medskip\noindent
In the rest of the proof of the lemma we assume
all underlying topological spaces are the
same, i.e., we have a diagram (\ref{equation-huge-2-ringed-topoi})
of sheaves on a site $\mathcal{C}$ (with $f_1 = f_2 = \text{id}$
on underlying topoi) and we are looking for
homomorphisms of sheaves of rings
$\mathcal{O}'_1 \to \mathcal{O}'_2$ fitting into it.
As ext groups we will use
$\Ext^k_{\mathcal{O}_1}(
\NL_{\mathcal{O}_1/\mathcal{O}_{\mathcal{B}_1}}, \mathcal{G}_2)$, $k = 0, 1$.

\medskip\noindent
Step 1. Construction of the obstruction class. Consider the sheaf
of sets
$$
\mathcal{E} = \mathcal{O}'_1 \times_{\mathcal{O}_2} \mathcal{O}'_2
$$
This comes with a surjective map $\alpha : \mathcal{E} \to \mathcal{O}_1$
and hence we can use $\NL(\alpha)$ instead of
$\NL_{\mathcal{O}_1/\mathcal{O}_{\mathcal{B}_1}}$, see
Modules on Sites, Lemma \ref{sites-modules-lemma-NL-up-to-qis}.
Set
$$
\mathcal{I}' =
\Ker(\mathcal{O}_{\mathcal{B}'_1}[\mathcal{E}] \to \mathcal{O}_1)
\quad\text{and}\quad
\mathcal{I} =
\Ker(\mathcal{O}_{\mathcal{B}_1}[\mathcal{E}] \to \mathcal{O}_1)
$$
There is a surjection $\mathcal{I}' \to \mathcal{I}$ whose kernel
is $\mathcal{J}_1\mathcal{O}_{\mathcal{B}'_1}[\mathcal{E}]$.
We obtain two homomorphisms of $\mathcal{O}_{\mathcal{B}'_2}$-algebras
$$
a : \mathcal{O}_{\mathcal{B}'_1}[\mathcal{E}] \to \mathcal{O}'_1
\quad\text{and}\quad
b : \mathcal{O}_{\mathcal{B}'_1}[\mathcal{E}] \to \mathcal{O}'_2
$$
which induce maps $a|_{\mathcal{I}'} : \mathcal{I}' \to \mathcal{G}_1$ and
$b|_{\mathcal{I}'} : \mathcal{I}' \to \mathcal{G}_2$. Both $a$ and $b$
annihilate $(\mathcal{I}')^2$. Moreover $a$ and $b$ agree on
$\mathcal{J}_1\mathcal{O}_{\mathcal{B}'_1}[\mathcal{E}]$
as maps into $\mathcal{G}_2$
because the left hand square of (\ref{equation-huge-2-ringed-topoi})
is commutative. Thus the difference
$b|_{\mathcal{I}'} - \nu \circ a|_{\mathcal{I}'}$
induces a well defined $\mathcal{O}_1$-linear map
$$
\xi : \mathcal{I}/\mathcal{I}^2 \longrightarrow \mathcal{G}_2
$$
which sends the class of a local section $f$ of $\mathcal{I}$ to
$a(f') - \nu(b(f'))$ where $f'$ is a lift of $f$ to a local
section of $\mathcal{I}'$. We let
$[\xi] \in \Ext^1_{\mathcal{O}_1}(\NL(\alpha), \mathcal{G}_2)$
be the image (see below).

\medskip\noindent
Step 2. Vanishing of $[\xi]$ is necessary. Let us write $\Omega =
\Omega_{\mathcal{O}_{\mathcal{B}_1}[\mathcal{E}]/\mathcal{O}_{\mathcal{B}_1}}
\otimes_{\mathcal{O}_{\mathcal{B}_1}[\mathcal{E}]} \mathcal{O}_1$.
Observe that $\NL(\alpha) = (\mathcal{I}/\mathcal{I}^2 \to \Omega)$
fits into a distinguished triangle
$$
\Omega[0] \to
\NL(\alpha) \to
\mathcal{I}/\mathcal{I}^2[1] \to
\Omega[1]
$$
Thus we see that $[\xi]$ is zero if and only if $\xi$
is a composition $\mathcal{I}/\mathcal{I}^2 \to \Omega \to \mathcal{G}_2$
for some map $\Omega \to \mathcal{G}_2$. Suppose there exists a
homomorphisms of sheaves of rings
$\varphi : \mathcal{O}'_1 \to \mathcal{O}'_2$ fitting into
(\ref{equation-huge-2-ringed-topoi}). In this case consider the map
$\mathcal{O}'_1[\mathcal{E}] \to \mathcal{G}_2$,
$f' \mapsto b(f') - \varphi(a(f'))$. A calculation
shows this annihilates $\mathcal{J}_1\mathcal{O}_{\mathcal{B}'_1}[\mathcal{E}]$
and induces a derivation
$\mathcal{O}_{\mathcal{B}_1}[\mathcal{E}] \to \mathcal{G}_2$.
The resulting linear map $\Omega \to \mathcal{G}_2$ witnesses the
fact that $[\xi] = 0$ in this case.

\medskip\noindent
Step 3. Vanishing of $[\xi]$ is sufficient. Let
$\theta : \Omega \to \mathcal{G}_2$ be a $\mathcal{O}_1$-linear map
such that $\xi$ is equal to
$\theta \circ (\mathcal{I}/\mathcal{I}^2 \to \Omega)$.
Then a calculation shows that
$$
b + \theta \circ d :
\mathcal{O}_{\mathcal{B}'_1}[\mathcal{E}]
\longrightarrow
\mathcal{O}'_2
$$
annihilates $\mathcal{I}'$ and hence defines a map
$\mathcal{O}'_1 \to \mathcal{O}'_2$ fitting into
(\ref{equation-huge-2-ringed-topoi}).

\medskip\noindent
Proof of (2) in the special case above. Omitted. Hint:
This is exactly the same as the proof of (2) of
Lemma \ref{lemma-huge-diagram}.
\end{proof}
